% Standalone candidate dossier for the posterior eigenfunction-defect calculus.
\documentclass[../main.tex]{subfiles}
\begin{document}

% === SOLUTION HEADER (proof provenance) =====================================
% ledger-node : lem:mm-posterior-defect
% refines     : lem:mm-posterior-defect
% bounded_by  : none (the ledger node carries no bounded_by edge)
% evidence_run: none
% checked_by  : agent
% author      : /root/prove_posterior_defect_par_06
% reviewer    : /root/review_posterior_defect_cold_08
% review      : research/reviews/2026-08-27-lem-mm-posterior-defect-proof-review.md
% date        : 2026-08-27
% ============================================================================

\section*{Solution: posterior eigenfunction-defect calculus}
\label{sec:sol-lem-mm-posterior-defect}

\paragraph{Refined statement.}
The regular setting in the following theorem makes explicit the domain and normalization
conventions implicit in Lemma~\ref{lem:mm-posterior-defect}.  For vectors $v,w\in\R^n$ we use
$(v\otimes w)_{ij}=v_iw_j$, and
$\operatorname{sym}M=(M+M^T)/2$.

\begin{theorem}[= Lemma~\ref{lem:mm-posterior-defect}]
\label{thm:sol-lem-mm-posterior-defect}
Let
\[
  \dd\mu(x)=Z^{-1}e^{-V(x)}\dd x
\]
be a smooth, strongly log-concave, centered isotropic probability measure on $\R^n$:
$V\in C^\infty(\R^n)$ and $\nabla^2V\succeq\varepsilon I_n$ for some $\varepsilon>0$.
Let $L=\Delta-\nabla V\cdot\nabla$ denote its Friedrichs self-adjoint realization in
$L^2(\mu)$, and let $f$ be a normalized first nonconstant eigenfunction,
\[
  -Lf=\lambda f,
  \qquad \E_\mu f=0,
  \qquad \E_\mu f^2=1.
\]
In particular, $f$ is understood in the generator domain; all identities below use the
closed Dirichlet-form interpretation, so there is no boundary condition hidden at infinity.

Take $X\sim\mu$ and an independent standard Brownian motion $B^{\mathrm{obs}}$, set
\[
  c_t=tX+B_t^{\mathrm{obs}},
  \qquad
  \mathcal F_t^{\mathrm{obs}}=\sigma(c_s:0\le s\le t),
\]
and let $\mu_t=\mathcal L(X\mid\mathcal F_t^{\mathrm{obs}})$.  Write $\E_t$ for integration
against $\mu_t$, $a_t=\E_tX$, and
\[
  m_t=\E_tf,
  \qquad
  g_t=\E_t[(f-m_t)(X-a_t)],
  \qquad
  H_t=\E_t[(f-m_t)(X-a_t)^{\otimes2}].
\]
For the posterior defect
\[
  R_t(x)=(c_t-tx)\cdot\nabla f(x),
  \qquad
  \rho_t=R_t-\E_tR_t,
\]
define, with the displayed tensor orientation,
\[
  b_t=\E_t\nabla f,
  \qquad
  C_t=\E_t[(X-a_t)\otimes\nabla f],
  \qquad
  u_t=\E_t[\rho_t(X-a_t)],
  \qquad
  K_t=\E_t[\rho_t(X-a_t)^{\otimes2}].
\]
Then, for every $t\ge0$, almost surely in the pointwise posterior identities,
\begin{align}
  \E_tR_t&=\lambda m_t,
  &\lambda g_t&=b_t+u_t,
  &\lambda H_t&=2\operatorname{sym}C_t+K_t,
  \label{eq:sol-mm-defect-identities}\\
  \E\Var_t(R_t)
  &=t\lambda-\lambda^2\int_0^t\E\abs{g_s}^2\dd s,
  &\E\int_0^t\norm{C_s}_{\HS}^2\dd s
  &\le \lambda-\lambda^2\abs{g_0}^2,
  \label{eq:sol-mm-defect-budgets}\\
  \E\E_t\abs{\nabla R_t}^2
  &\le t\lambda^2+t^2\lambda.
  \label{eq:sol-mm-defect-gradient}
\end{align}
Here the un-subscripted expectation averages the planted observation model, and
$\Var_t(R_t)=\E_t\rho_t^2$.
\end{theorem}

\begin{proof}
\noindent\emph{Posterior and innovation equations.}
The likelihood of the observation path up to time $t$, conditional on $X=x$, depends on that
path only through $c_t$ and is proportional to
$\exp(c_t\cdot x-t\abs{x}^2/2)$.  Hence
\begin{equation}
  \dd\mu_t(x)
  =\frac{\exp(c_t\cdot x-t\abs{x}^2/2)}
  {\int\exp(c_t\cdot y-t\abs{y}^2/2)\dd\mu(y)}\,\dd\mu(x).
  \label{eq:sol-mm-posterior-density}
\end{equation}
Put
\[
  W_t=c_t-\int_0^ta_s\dd s.
\]
Since $\E[X\mid\mathcal F_t^{\mathrm{obs}}]=a_t$, the process $W$ is a continuous
$\mathcal F_t^{\mathrm{obs}}$-local martingale with quadratic covariation
$[W^i,W^j]_t=\delta_{ij}t$.  L\'evy's characterization therefore makes $W$ an
$n$-dimensional Brownian motion.  Applying It\^o's formula to
\eqref{eq:sol-mm-posterior-density} gives
\begin{equation}
  \dd\mu_t(x)=\mu_t(x)(x-a_t)\cdot\dd W_t.
  \label{eq:sol-mm-density-sde}
\end{equation}
Consequently, for every fixed scalar $\phi\in L^2(\mu)$,
\begin{equation}
  \dd\E_t\phi=\Cov_t(\phi,X)\cdot\dd W_t.
  \label{eq:sol-mm-filtering}
\end{equation}

We record the localization and domain justification used throughout.  First take bounded smooth
truncations of $\phi$ and stop at the first time when the stochastic-integral quadratic
variation or the posterior moments occurring in the calculation reach level $N$.  Equations
\eqref{eq:sol-mm-density-sde}--\eqref{eq:sol-mm-filtering} are then ordinary bounded It\^o
identities.  Conditional expectation is a contraction in $L^2$, so bounded truncations of $f$
and of each component of $\nabla f$ converge to their square-integrable conditional-expectation
martingales.  Their stopped stochastic integrals therefore converge in $L^2$, and the stopping
times increase to infinity.  This removes the stopping without losing any equality below.

For the spatial integrations by parts, the posterior potential is
\[
  V_t(x)=V(x)-c_t\cdot x+\frac t2\abs{x}^2,
  \qquad \nabla^2V_t\succeq(\varepsilon+t)I_n.
\]
Thus $\mu_t$ has all polynomial moments.  The closed Dirichlet form
\[
  \mathcal E_t(h,k)=\E_t[\nabla h\cdot\nabla k],
  \qquad h,k\in W^{1,2}(\mu_t),
\]
has generator
\begin{equation}
  L_t=\Delta-\nabla V_t\cdot\nabla
      =L+(c_t-tx)\cdot\nabla.
  \label{eq:sol-mm-Lt}
\end{equation}
The identities below are first paired with compactly supported cutoffs of the coordinate and
centered-quadratic tests.  The Gaussian tails of $\mu_t$, the $L^2$ bounds established by the
planted coupling below, and closedness of $\mathcal E_t$ let the cutoff radius tend to infinity.
Equivalently, the coordinates and centered quadratics belong to the required local form domain,
and $f$ belongs to the maximal generator domain of $L_t$ because both $f$ and $L_tf$ are in
$L^2(\mu_t)$.  This specifies the domain in every use of posterior integration by parts.

\medskip
\noindent\emph{The defect identities.}
The eigenfunction equation and \eqref{eq:sol-mm-Lt} give the exact generator relation
\begin{equation}
  -L_tf=\lambda f-R_t.
  \label{eq:sol-mm-generator-defect}
\end{equation}
Pairing this identity with the constant function in the posterior Dirichlet form yields
\[
  0=\E_t[-L_tf]=\lambda m_t-\E_tR_t,
\]
which is the first identity in \eqref{eq:sol-mm-defect-identities}.

Pair next with the coordinate $x_i-a_{t,i}$.  Posterior integration by parts gives
\[
  \E_t[(-L_tf)(X_i-a_{t,i})]=\E_t\partial_i f=(b_t)_i.
\]
Because $\E_t(X-a_t)=0$ and $\rho_t=R_t-\lambda m_t$, the same left-hand side, evaluated with
\eqref{eq:sol-mm-generator-defect}, is $(\lambda g_t-u_t)_i$.  Thus
$\lambda g_t=b_t+u_t$.

Finally fix an arbitrary symmetric matrix $D$ and set
\[
  Q_D(x)=(x-a_t)^TD(x-a_t)-\Tr(DA_t).
\]
Then $\E_tQ_D=0$, $\nabla Q_D=2D(x-a_t)$, and our orientation for $C_t$ gives
\begin{align*}
  \E_t[(-L_tf)Q_D]
  &=\E_t[\nabla f\cdot\nabla Q_D]
    =2\langle D,\operatorname{sym}C_t\rangle_{\HS},\\
  \E_t[(\lambda f-R_t)Q_D]
  &=\langle D,\lambda H_t-K_t\rangle_{\HS}.
\end{align*}
The two expressions agree by \eqref{eq:sol-mm-generator-defect}.  Since they agree for every
symmetric $D$ and both $H_t,K_t$ are symmetric, we obtain
$\lambda H_t=2\operatorname{sym}C_t+K_t$.

\medskip
\noindent\emph{The centered-defect budget.}
At the planted point, the definition of the observation gives the cancellation
\begin{equation}
  R_t(X)=B_t^{\mathrm{obs}}\cdot\nabla f(X).
  \label{eq:sol-mm-planted-defect}
\end{equation}
The usual energy identity for the normalized eigenfunction is
\begin{equation}
  \E_\mu\abs{\nabla f}^2
  =\langle f,-Lf\rangle_{L^2(\mu)}=\lambda.
  \label{eq:sol-mm-gradient-energy}
\end{equation}
Since $B_t^{\mathrm{obs}}$ is independent of $X$, has mean zero, and has covariance $tI_n$,
the conditional-law property and \eqref{eq:sol-mm-planted-defect} imply
\begin{equation}
  \E\E_tR_t^2
  =\E\abs{B_t^{\mathrm{obs}}\cdot\nabla f(X)}^2
  =t\lambda.
  \label{eq:sol-mm-defect-second-moment}
\end{equation}
In particular, the $L^2(\mu_t)$ assertion used in the domain paragraph holds for almost every
observation.

Apply \eqref{eq:sol-mm-filtering} to $f$.  Since $m_0=0$ and
$\Cov_t(f,X)=g_t$, the square-integrable martingale $m$ satisfies
\[
  \dd m_t=g_t\cdot\dd W_t,
  \qquad
  \E m_t^2=\int_0^t\E\abs{g_s}^2\dd s.
\]
The second equality is first the stopped It\^o isometry; it remains an equality after removal
because $m_t=\E[f(X)\mid\mathcal F_t^{\mathrm{obs}}]$ is bounded in $L^2$.
Using $\E_tR_t=\lambda m_t$ and conditional variance decomposition in
\eqref{eq:sol-mm-defect-second-moment} now gives
\[
  \E\Var_t(R_t)
  =\E\E_tR_t^2-\E(\E_tR_t)^2
  =t\lambda-\lambda^2\int_0^t\E\abs{g_s}^2\dd s.
\]

\medskip
\noindent\emph{The $C_t$ martingale budget.}
Apply the vector form of \eqref{eq:sol-mm-filtering} to the fixed vector field $\nabla f$.
With the orientation specified in the theorem,
\[
  \dd b_t=C_t^T\dd W_t.
\]
The martingale $b_t=\E[\nabla f(X)\mid\mathcal F_t^{\mathrm{obs}}]$ is square-integrable by
\eqref{eq:sol-mm-gradient-energy}.  The stopped vector It\^o isometry and the same $L^2$
stopping-removal argument therefore give
\begin{equation}
  \E\int_0^t\norm{C_s}_{\HS}^2\dd s
  =\E\abs{b_t}^2-\abs{b_0}^2.
  \label{eq:sol-mm-C-isometry}
\end{equation}
Conditional Jensen and the tower property show that
$\E\abs{b_t}^2\le\E\abs{\nabla f(X)}^2=\lambda$.  At time zero,
$c_0=0$ and hence $R_0=\rho_0=u_0=0$; the already proved coordinate identity gives
$b_0=\lambda g_0$.  Substitution in \eqref{eq:sol-mm-C-isometry} proves the second estimate in
\eqref{eq:sol-mm-defect-budgets}.

\medskip
\noindent\emph{The averaged gradient-defect budget.}
Differentiating $R_t$ with respect to its spatial argument gives
\begin{equation}
  \nabla R_t(x)=\nabla^2f(x)(c_t-tx)-t\nabla f(x).
  \label{eq:sol-mm-gradient-defect-formula}
\end{equation}
Evaluate at the planted point and use $c_t-tX=B_t^{\mathrm{obs}}$.  Independence and centering
of $B_t^{\mathrm{obs}}$ make the cross term vanish, while its covariance is $tI_n$.  Therefore
\begin{equation}
  \E\E_t\abs{\nabla R_t}^2
  =t\E_\mu\norm{\nabla^2f}_{\HS}^2+t^2\E_\mu\abs{\nabla f}^2.
  \label{eq:sol-mm-gradient-defect-exact}
\end{equation}
For the Friedrichs realization, the integrated weighted Bochner identity, justified by the same
compact cutoff approximation, is
\[
  \E_\mu(Lf)^2
  =\E_\mu\norm{\nabla^2f}_{\HS}^2
   +\E_\mu\langle\nabla^2V\nabla f,\nabla f\rangle.
\]
Convexity of $V$, the eigenfunction equation, and normalization consequently imply
\[
  \E_\mu\norm{\nabla^2f}_{\HS}^2
  \le\E_\mu(Lf)^2=\lambda^2.
\]
Combining this with \eqref{eq:sol-mm-gradient-energy} and
\eqref{eq:sol-mm-gradient-defect-exact} gives \eqref{eq:sol-mm-defect-gradient}.
All limiting operations used here are either $L^2$ limits of the two conditional-expectation
martingales or closed-form $L^2$ limits of the spatial cutoff approximations, so the auxiliary
stopping and spatial cutoffs have now been removed completely.
\end{proof}

\begin{remark}[Excluded conditional consequence]
No use is made of the imported, unreviewed node \texttt{thm:letwin-qcts}.  In particular, a
whitened quadratic estimate for $K_t$ conditional on that node is outside this theorem and this
dossier.  The identities proved here do not unwhiten the high-covariance part of $H_t$, do not
establish the high-incidence occupation estimate, and supply no approximation passage to an
arbitrary log-concave law.
\end{remark}

\paragraph{Obstructions respected.}
The ledger node has no formal \texttt{bounded\_by} edge.  The route fences were nevertheless
checked individually.  The proof uses no cut or slice estimate, so \texttt{obs:two-tail},
\texttt{obs:circularity}, and \texttt{obs:rank-one-refuted} are not engaged.  It derives exact
coordinate and full symmetric-matrix integration-by-parts identities rather than a
dimension-free quadratic-chaos theorem from radial or projection tests, so it does not cross
\texttt{obs:proj-ceiling}.  It uses neither the crude covariance integral nor a relative
covariance occupation bound, respecting \texttt{obs:crude-insufficient} and
\texttt{obs:relative-ceiling}.  Finally, it never bounds a posterior covariance operator norm,
never unwhitens a tensor, and never invokes the refuted truncated-exponential variable-weight
Stein shortcut; hence it also respects the covariance-spike and spectral-unwhitening warnings.
The transport and needle warnings are inapplicable because neither mechanism occurs.

\end{document}
